\section{Experiments}
\label{sec:experiments}

% Subsections still to be written by their owners:
%   - SAE training and evaluation
%   - Layer-wise probes and baselines
%   - Causal intervention design

\subsection{From free-form feedback to candidate presentation}
\label{sec:exp-protocol}

The first collection used a free-form protocol: the Solver revised its
answer after reading the Critic's full feedback text. The same questions and
models were then collected again with the candidate presentation, in which
the Solver chooses between its first answer and the Critic's recommendation,
shown as two labeled options. We compare the Adoption share of the binary
cohort under the two protocols. Because the earlier collection was labeled
with an earlier rule set, we report it with both its own rules and the
LLM-derived labels.

\subsection{Candidate-order control}
\label{sec:exp-order}

In the candidate presentation, the Solver sees its own first answer and the
Critic's recommendation as two labeled options. Each question is assigned
one of two orders before generation, either A1 first or Critic first. The
assignment is seeded and balanced within each hop group, so it is
independent of the question's content and outcome. We compare the Adoption
share of the binary cohort between the two orders in every partition, with
95\% Wilson intervals. A difference would show that answer position, rather
than the content of the feedback, influences the Solver's choice.

\subsection{Synthetic-question pilot}
\label{sec:exp-synthetic}

To enlarge the Adoption and Retention cohort, we generated synthetic
multi-hop questions in the MuSiQue format with Claude Opus~5. Each question
comes with its supporting paragraphs, three generated distractor
paragraphs, and filler paragraphs drawn from MuSiQue. Before adding them to
the training data, we ran a pilot to check whether they produce feedback
behavior comparable to MuSiQue.

The pilot collected 100 synthetic questions with the same models, protocol,
candidate presentation, and decoding settings as the MuSiQue collection. As
a control, we took 100 MuSiQue questions from the frozen training partition
(seed 20261030), matched to the pilot's hop mix: 34 two-hop, 20 \textit{3hop1},
13 \textit{3hop2}, 10 \textit{4hop1}, 9 \textit{4hop2}, and 14 \textit{4hop3}
questions. Both groups were labeled with the same LLM labeler and settings
(Section~\ref{sec:method-labels}). We also measured where the two groups
differ: A1 and Critic accuracy, whether a lexical retriever ranks the
supporting paragraphs above the rest, whether wrong answers come from a
non-supporting paragraph, and how often the Critic refuses.

\subsection{Paired filler ablation}
\label{sec:exp-fillers}

The pilot suggested that the filler paragraphs, drawn at random from
MuSiQue, rarely compete with the evidence. To test this directly, we rebuilt
the same 100 synthetic questions with only the filler paragraphs replaced.
Each filler is swapped, in place, for a paragraph drawn (with a fixed seed)
from the 20 eligible MuSiQue paragraphs most similar to the question and its
supporting paragraphs (TF-IDF). Every
question, answer, alias, decomposition, supporting paragraph, and generated
distractor is unchanged. A filler that mentions any hop's answer or alias is
excluded, so a replacement cannot restate or contradict the evidence.

Both conditions are collected in the same environment, in the original
question order, so per-question seeds and candidate orders are identical
and the fillers are the only difference. The comparison is paired: for each
question, we record whether its label changes between the original and the
related fillers.

We repeated the comparison once on 100 further synthetic questions from the
same generation run, sampled with a separate seed and not inspected before
collection, with the analysis fixed in advance.
